%!TeX encoding = UTF-8
%!TeX program = pdflatex
% EMS Scalability Enhancement Proposal - LaTeX Template
% Simple, clean template for quick professor approval
% Compatible with WinEdt 11

\documentclass[12pt,a4paper,oneside]{report}

% ============================================================================
% PACKAGES
% ============================================================================

\usepackage[utf8]{inputenc}
\usepackage[T1]{fontenc}
\usepackage[margin=1in]{geometry}
\usepackage{setspace}
\usepackage{amsmath}
\usepackage{amssymb}
\usepackage{graphicx}
\usepackage{booktabs}
\usepackage{multirow}
\usepackage{array}
\usepackage{xcolor}
\usepackage{hyperref}
\usepackage{fancyhdr}
\usepackage{float}
\usepackage{caption}
\usepackage{subcaption}
\usepackage{listings}
\usepackage{xspace}
\usepackage{enumerate}
\usepackage{url}
\usepackage{natbib}
\usepackage{tocloft}
\usepackage{tikz}
\usepackage{pdflscape}
\usepackage{enumitem}
\usepackage{threeparttable}
\usepackage{tocbibind}
\usepackage{calc}

% ============================================================================
% CONFIGURATION
% ============================================================================

% Hyperlink colors (blue links, no boxes)
\hypersetup{
    colorlinks=true,
    linkcolor=blue,
    filecolor=blue,
    urlcolor=blue,
    citecolor=blue
}

% Pagestyle
\pagestyle{fancy}
\fancyhf{}
\fancyhead[R]{\thepage}
\fancyhead[L]{EMS Scalability Enhancement}
\fancyfoot[C]{\thepage}

\onehalfspacing

\setcounter{secnumdepth}{3}
\tocbibind

% Define custom colors
\definecolor{gray10}{gray}{10}

% Code listing style (for SQL, Java, etc.)
\lstset{
    basicstyle=\small\ttfamily,
    breaklines=true,
    breakatwhitespace=true,
    frame=single,
    numbers=left,
    numberstyle=\tiny,
    keywordstyle=\color{blue},
    commentstyle=\color{gray},
    stringstyle=\color{red},
    language=Java,
    showstringspaces=false
}

% ============================================================================
% CUSTOM COMMANDS
% ============================================================================

% Highlight box for key metrics
\newcommand{\metricbox}[2]{%
    \fcolorbox{black}{gray10}{%
        \parbox{0.45\textwidth}{%
            \textbf{#1:} #2%
        }%
    }%
}

% ============================================================================
% DOCUMENT METADATA
% ============================================================================

\title{\LARGE\textbf{EMS Scalability Enhancement Proposal:\\
Week 2 Backend Optimization Plan}}

\author{
    \Large [Your Name] \\
    \large Sophomore, Computer Science \\
    \large [University Name] \\
    \large Email: [your.email@university.edu] \\
    \large Student ID: [Your Student ID]
}

\date{\Large February 13, 2026}

% ============================================================================
% DOCUMENT CONTENT
% ============================================================================

\begin{document}

\maketitle

% ============================================================================
% ABSTRACT
% ============================================================================

\begin{abstract}

\noindent\textbf{Problem Statement:} The EMS backend (Spring Boot 3.3.5 + MySQL 8.0) currently handles \textbf{100-150 concurrent users} but fails critically under realistic university load. Baseline load testing revealed the system cannot support peak registration periods, where 500-2,500 concurrent students attempt to register for popular events.

\vspace{0.5em}
\noindent
\metricbox{Critical Gap}{1.8\% success rate under high contention}
\hfill
\metricbox{Target}{98-100\% success rate after optimization}

\noindent\textbf{Root Cause Analysis:} Systematic load testing using k6 identified \textbf{five specific bottlenecks}:

\begin{table}[H]
\centering
\small
\begin{tabular}{@{}clll@{}}
\toprule
\textbf{Priority} & \textbf{Bottleneck} & \textbf{Current State} & \textbf{Impact} \\
\midrule
\textbf{P1 - CRITICAL} & Connection Pool & 10 connections & \textcolor{red}{82.8\% timeout errors} \\
\textbf{P2 - HIGH} & JWT DB Queries & Query on every request & 46,232 redundant queries \\
\textbf{P3 - MEDIUM} & Missing Indexes & Full table scans & 10-50ms wasted per query \\
\textbf{P4 - MEDIUM} & Fixed Backoff & 50ms fixed delay & Retry collisions \\
\textbf{P5 - LOW} & Synchronous Email & Blocks threads & 300ms added latency \\
\bottomrule
\end{tabular}
\caption{Identified Bottlenecks (Prioritized by Impact)}
\end{table}

\noindent\textbf{Key Success:} Optimistic locking with \texttt{@Version} annotation \textbf{prevented all capacity violations} (zero overbooking across 2,766 registrations tested). The concurrency control mechanism is production-ready.

\noindent\textbf{Proposed Solution:} Implement \textbf{five sequential optimizations} over 7 days (Week 2: Days 8-14) to improve system capacity from 150 to \textbf{650-750 concurrent users} (5-7x improvement).

\begin{table}[H]
\centering
\small
\begin{tabular}{@{}llll@{}}
\toprule
\textbf{Optimization} & \textbf{Time} & \textbf{Capacity Gain} & \textbf{Technical Change} \\
\midrule
P1: Connection Pool & 1 hour & +300 users & \texttt{max-pool-size: 10 to 50} \\
P2: JWT Optimization & 4 hours & +150 users & Add userId/role to token \\
P3: Database Indexes & 2 hours & +75 users & 4 CREATE INDEX statements \\
P4: Exponential Backoff & 2 hours & +40 users & Retry logic with jitter \\
P5: Async Email & 3 hours & +35 users & \texttt{@Async} + thread pool \\
\midrule
\textbf{Total} & \textbf{12 hours} & \textbf{+500-600 users} & \textbf{5-7x improvement} \\
\bottomrule
\end{tabular}
\caption{Week 2 Optimization Plan with Expected Outcomes}
\end{table}

\noindent\textbf{Expected Week 2 Results:}
\begin{itemize}[leftmargin=1.5em,topsep=0.3em,itemsep=0.2em]
    \item Scenario 1 success rate: 1.8\% to \textbf{98-100\%} (54x improvement)
    \item Scenario 2 success rate: 12.5\% to \textbf{75-85\%} (6-7x improvement)
    \item Average response time: 9.15s to \textbf{less than 200ms} (45x improvement)
    \item HTTP failure rate: 98\% to \textbf{less than 2\%} (49x reduction)
\end{itemize}

\noindent\textbf{Request for Approval:}

\textbf{Seeking permission to:}
\begin{enumerate}[leftmargin=1.5em,topsep=0.3em,itemsep=0.2em]
    \item Proceed with Week 2 optimizations (Days 8-14)
    \item Modify backend configuration (\texttt{application.properties})
    \item Refactor JWT authentication logic (add claims to token)
    \item Add database indexes to production schema
    \item Re-run load tests to validate improvements
\end{enumerate}

\textbf{No architectural changes required} -- all optimizations are configuration updates and code-level improvements to existing components.

\end{abstract}

\clearpage

% ============================================================================
% TABLE OF CONTENTS
% ============================================================================

\tableofcontents

\clearpage

% ============================================================================
% CHAPTER 1: INTRODUCTION
% ============================================================================

\chapter{Introduction}

\section{Project Context}

The \textbf{Extracurricular Management System (EMS)} is a university event management platform that automates the complete event lifecycle:

\begin{itemize}[leftmargin=1.5em,topsep=0.3em,itemsep=0.2em]
    \item \textbf{Event Organizers} create event proposals
    \item \textbf{Administrators} review and approve proposals
    \item \textbf{Students} register for approved events with real-time capacity enforcement
\end{itemize}

\textbf{Current Project Status:}
\begin{itemize}[leftmargin=1.5em,topsep=0.3em,itemsep=0.2em]
    \item \textbf{Phase:} OOAD design complete (90\%), transitioning to production implementation
    \item \textbf{Technical Stack:} Spring Boot 3.3.5 (Java 21) + React 19 + MySQL 8.0
    \item \textbf{Core Features:} JWT authentication, role-based access control, optimistic locking
    \item \textbf{Documentation:} SRS v5 (92/100), complete use case diagrams, sequence diagrams
\end{itemize}

\section{Motivation for Scalability Enhancement}

\textbf{Academic Context:}

As part of the OOAD final project, demonstrating production readiness requires validating that the system can handle realistic university workloads. Load testing is a critical phase of software engineering to identify and resolve performance bottlenecks before deployment.

\vspace{0.5em}
\textbf{Real-World Need:}

Universities experience \textbf{peak registration periods} where hundreds of students simultaneously attempt to register for popular events (e.g., semester start, high-demand guest speakers, limited-capacity workshops). The current system capacity (100-150 users) is insufficient for a mid-sized university (5,000-10,000 students).

\vspace{0.5em}
\textbf{Technical Challenge:}

The system must handle:
\begin{itemize}[leftmargin=1.5em,topsep=0.3em,itemsep=0.2em]
    \item \textbf{Flash sale scenarios:} 500 students competing for 50 spots in a popular event
    \item \textbf{Distributed load:} 500 concurrent students registering across multiple events
    \item \textbf{Concurrency correctness:} Zero capacity violations (no overbooking)
    \item \textbf{Response time:} Less than 500ms for acceptable user experience
\end{itemize}

\section{Objectives}

\textbf{Primary Objective:}

Scale the EMS backend to handle \textbf{2,500 concurrent users} during peak registration periods while maintaining data integrity (no capacity violations) and acceptable response times (less than 500ms).

\vspace{0.5em}
\textbf{Week 2 Objective (This Proposal):}

Achieve \textbf{650-750 concurrent users} (5-7x improvement) through backend optimizations:

\begin{enumerate}[leftmargin=1.5em,topsep=0.3em,itemsep=0.2em]
    \item Expand HikariCP connection pool (10 to 50 connections)
    \item Optimize JWT authentication (eliminate database queries)
    \item Add database indexes (speed up lookups)
    \item Implement exponential backoff (reduce retry collisions)
    \item Make email notifications asynchronous (free request threads)
\end{enumerate}

\textbf{Success Criteria:}
\begin{itemize}[leftmargin=1.5em,topsep=0.3em,itemsep=0.2em]
    \item Scenario 1 (single event) success rate greater than 95\%
    \item Scenario 2 (distributed load) success rate greater than 85\%
    \item Average response time less than 200ms
    \item HTTP failure rate less than 2\%
    \item Zero capacity violations maintained (optimistic locking continues to work)
\end{itemize}

\section{Methodology}

\textbf{Approach:} Systematic, data-driven performance engineering following industry best practices:

\begin{enumerate}[leftmargin=1.5em,topsep=0.3em,itemsep=0.3em]
    \item \textbf{Measure (Week 1 - Complete):}
    \begin{itemize}[topsep=0.2em,itemsep=0.1em]
        \item Establish baseline performance metrics
        \item Identify bottlenecks through load testing
        \item Quantify impact of each bottleneck
    \end{itemize}

    \item \textbf{Optimize (Week 2 - This Proposal):}
    \begin{itemize}[topsep=0.2em,itemsep=0.1em]
        \item Prioritize fixes by ROI (impact / effort)
        \item Implement optimizations sequentially
        \item Validate each optimization independently
    \end{itemize}

    \item \textbf{Validate (Week 2 - Days 12-14):}
    \begin{itemize}[topsep=0.2em,itemsep=0.1em]
        \item Re-run identical load tests
        \item Compare before/after metrics
        \item Document improvements and lessons learned
    \end{itemize}

    \item \textbf{Scale (Week 3-4 - Future):}
    \begin{itemize}[topsep=0.2em,itemsep=0.1em]
        \item Horizontal scaling (multiple backend instances)
        \item Advanced caching (Redis)
        \item Achieve final 2,500-user target
    \end{itemize}
\end{enumerate}

\textbf{Load Testing Tool:} k6 v1.5.0 (industry-standard, open-source)

\textbf{Test Environment:} Localhost (eliminates network latency as variable, focuses on application bottlenecks)

\textbf{Test Scenarios:}
\begin{itemize}[leftmargin=1.5em,topsep=0.3em,itemsep=0.2em]
    \item \textbf{Scenario 1 (Extreme Contention):} 500 VUs (virtual users) simultaneously register for 1 event with 50 capacity
    \item \textbf{Scenario 2 (Distributed Load):} 0 to 500 VUs over 5 minutes, registering across 50 events
\end{itemize}

\section{Scope}

\textbf{In Scope:}
\begin{itemize}[leftmargin=1.5em,topsep=0.3em,itemsep=0.2em]
    \item Backend performance optimizations (Spring Boot + MySQL)
    \item HikariCP connection pool configuration
    \item JWT token structure changes
    \item Database schema indexing
    \item Retry logic improvements
    \item Email service asynchronization
\end{itemize}

\textbf{Out of Scope:}
\begin{itemize}[leftmargin=1.5em,topsep=0.3em,itemsep=0.2em]
    \item Frontend performance optimization (React rendering, bundle size)
    \item Network optimization (CDN, load balancing)
    \item Database architecture changes (sharding, replication)
    \item Cloud deployment (addressed in Week 3-4)
    \item Security enhancements (separate audit planned)
\end{itemize}

\section{Document Structure}

This proposal consists of three complementary documents:

\begin{enumerate}[leftmargin=1.5em,topsep=0.3em,itemsep=0.3em]
    \item \textbf{BASELINE\_SUMMARY.md} (Quick Reference)
    \begin{itemize}[topsep=0.2em,itemsep=0.1em]
        \item 1-page overview of key findings
        \item Metrics tables for quick scanning
        \item Week 2 optimization checklist
    \end{itemize}

    \item \textbf{BASELINE\_RESULTS.md} (Detailed Technical Report)
    \begin{itemize}[topsep=0.2em,itemsep=0.1em]
        \item Complete test configuration and execution details
        \item Scenario 1 and 2 results with analysis
        \item Infrastructure metrics (HikariCP, MySQL, JVM)
        \item Optimistic locking validation
        \item Code locations and SQL queries
    \end{itemize}

    \item \textbf{BASELINE\_SUGGESTIONS.md} (Implementation Roadmap)
    \begin{itemize}[topsep=0.2em,itemsep=0.1em]
        \item Detailed bottleneck analysis with queuing theory
        \item Priority recommendations (P1-P5) with ROI calculations
        \item Week 2 day-by-day implementation plan
        \item Risk assessment and mitigation strategies
        \item Verification strategy for post-optimization testing
    \end{itemize}
\end{enumerate}

\textbf{This document} provides the executive overview and introduction for quick professor review and approval.

\section{Why This Matters}

\textbf{Engineering Maturity:}

This work demonstrates application of software engineering principles learned in OOAD:
\begin{itemize}[leftmargin=1.5em,topsep=0.3em,itemsep=0.2em]
    \item \textbf{Separation of Concerns:} Load testing identified layer-specific bottlenecks
    \item \textbf{Design Patterns:} Optimistic locking validated under stress
    \item \textbf{Performance Engineering:} Systematic measurement to optimization cycle
\end{itemize}

\textbf{Production Readiness:}

Real-world systems must handle peak load. This project goes beyond ``it works on my laptop'' to ``it works under university-scale concurrent load.''

\textbf{Academic Contribution:}

Applying queuing theory to model connection pool behavior, calculating theoretical maximum throughput, and validating predictions with empirical testing demonstrates research rigor.

\section{Key Terminology}

\begin{table}[H]
\centering
\small
\begin{tabular}{@{}ll@{}}
\toprule
\textbf{Term} & \textbf{Definition} \\
\midrule
VU (Virtual User) & Simulated user in k6 load testing \\
Concurrent Users & Users actively making requests simultaneously \\
Connection Pool & Pre-allocated database connections (HikariCP) \\
Optimistic Locking & Concurrency control using version numbers \\
JWT & JSON Web Token (stateless authentication) \\
p95 Response Time & 95th percentile response time \\
Throughput & Requests processed per second \\
Capacity & Max concurrent users with acceptable performance \\
\bottomrule
\end{tabular}
\caption{Key Technical Terminology}
\end{table}

\clearpage

% ============================================================================
% CHAPTER 2: NEXT STEPS
% ============================================================================

\chapter{Next Steps}

\section{Immediate (Pending Approval)}

\begin{enumerate}[leftmargin=1.5em,topsep=0.3em,itemsep=0.2em]
    \item Professor reviews this proposal and supporting documents
    \item Approval to proceed with Week 2 optimizations (12 hours of work)
    \item Begin implementation on Day 8
\end{enumerate}

\section{Week 2 Timeline (7 days)}

\begin{itemize}[leftmargin=1.5em,topsep=0.3em,itemsep=0.2em]
    \item \textbf{Days 8-11:} Implement P1-P5 optimizations (1 per day, plus async email)
    \item \textbf{Days 12-13:} Re-run load tests, capture new metrics
    \item \textbf{Day 14:} Document results, compare to baseline, prepare Week 3 plan
\end{itemize}

\section{Deliverable}

\begin{itemize}[leftmargin=1.5em,topsep=0.3em,itemsep=0.2em]
    \item Week 2 Results Report showing before/after comparison
    \item Updated system capacity: 150 to 650-750 concurrent users
    \item Recommendation for Week 3 (horizontal scaling) or further optimization
\end{itemize}

\vspace{2em}

\noindent
\textbf{Respectfully submitted for approval,} \\
\\[1em]
\noindent
[Your Name] \\
Sophomore, Computer Science \\
[University Name] \\
[Date]

% ============================================================================
% APPENDIX
% ============================================================================

\clearpage

\appendix

\chapter{Supporting Documentation}

The following documents provide detailed technical analysis and are attached as appendices:

\begin{enumerate}[leftmargin=1.5em,topsep=0.3em,itemsep=0.3em]
    \item \textbf{BASELINE\_SUMMARY.md}

    Quick reference showing key metrics and optimization checklist

    \item \textbf{BASELINE\_RESULTS.md}

    Comprehensive test results including:
    \begin{itemize}[topsep=0.2em,itemsep=0.1em]
        \item Scenario 1: Single Event Stress Test (500 VUs to 1 event)
        \item Scenario 2: Distributed Load Test (500 VUs to 50 events)
        \item Infrastructure metrics validation
        \item Optimistic locking validation (zero capacity violations)
    \end{itemize}

    \item \textbf{BASELINE\_SUGGESTIONS.md}

    Implementation roadmap including:
    \begin{itemize}[topsep=0.2em,itemsep=0.1em]
        \item Bottleneck deep dive with queuing theory analysis
        \item Priority recommendations (P1-P5) with ROI calculations
        \item Day-by-day implementation schedule
        \item Risk assessment and mitigation strategies
    \end{itemize}
\end{enumerate}

\chapter{Test Data Files}

Available upon request:
\begin{itemize}[leftmargin=1.5em,topsep=0.3em,itemsep=0.2em]
    \item \texttt{scenario1-single-event.json} - k6 metrics (Scenario 1)
    \item \texttt{scenario1-console.txt} - k6 terminal output (179KB)
    \item \texttt{scenario2-distributed.json} - k6 metrics (Scenario 2)
    \item \texttt{scenario2-console.txt} - k6 terminal output (142KB)
    \item \texttt{student\_tokens.csv} - 10,000 pre-generated JWT tokens
\end{itemize}

% ============================================================================
% END DOCUMENT
% ============================================================================

\end{document}
